% from ega4-18-fr.tex
\oldpage[IV-4]{113}

l'homomorphisme $\operatorname{astr}_{\sh{B}|U/\sh{O}_U}$ soit bijectif. Le fait que $b)$ implique $a)$ résulte alors de (18.2.3).

\begin{corollary}[18.2.5]
\label{IV.18.2.5}
Soient $Y$ un préschéma quasi-compact ou localement noethérien, $f:X\to Y$ un morphisme fini et de présentation finie; posons $\sh{B}=f_*(\sh{O}_X)$. Supposons que pour tout point fermé $y$ de $Y$, $\sh{B}_y$ soit un $\sh{O}_{Y,y}$-module libre et $\sh{B}_y\otimes_{\sh{O}_{Y,y}}\kappa(y)$ une $\kappa(y)$-algèbre séparable. Alors $f$ est étale.
\end{corollary}

En effet, il résulte de (18.2.4) que tout point fermé $y$ de $Y$ possède un voisinage ouvert $U$ tel que la restriction $f^{-1}(U)\to U$ de $f$ soit étale; la conclusion résulte de ce que dans les deux cas considérés, toute partie fermée non vide de $Y$ contient un point fermé (5.1.11 et $0_{\mathrm I}$, 2.1.3).

\begin{corollary}[18.2.6]
\label{IV.18.2.6}
Si $f:X\to Y$ est un morphisme fini et étale, et si l'on pose $\sh{B}=f_*(\sh{O}_X)$, alors le $\sh{O}_Y$-homomorphisme $\operatorname{Tr}_{\sh{B}|\sh{O}_Y}:\sh{B}\to\sh{O}_Y$ (18.2.2) (aussi noté $\operatorname{Tr}_f$) est surjectif.
\end{corollary}

La question étant locale, on peut, en vertu de (18.2.3) supposer que $Y=\Spec(A)$, $X=\Spec(B)$ avec $\sh{B}=\widetilde B$, $B$ étant un $A$-module libre; comme en vertu de (18.2.3) la forme bilinéaire (18.2.1.1) est non dégénérée, cela entraîne en particulier que la forme linéaire $\operatorname{Tr}_{B/A}$ est surjective.

\begin{remarks}[18.2.7]
\EGAReaderTarget{ega4.c04.item.18-2-7}{IV.18.2.7}
\begin{enumerate}
  \item[(i)] Lorsque $f:X\to Y$ est un morphisme fini tel que $f_*(\sh{O}_X)$ soit un $\sh{O}_Y$-Module localement libre (resp. localement libre de rang $n$), on dit encore que $f$ est un morphisme \emph{fini localement libre} (resp. \emph{localement libre de rang $n$}). Cette condition, en vertu de (18.2.3), est vérifiée si $f$ est un morphisme fini étale, mais n'implique pas à elle seule que $f$ soit étale, comme le montre l'exemple où $X=\Spec(K)$ et $Y=\Spec(k)$ sont des spectres de corps, $K$ étant une extension finie non séparable de $k$. Lorsque $f:X\to Y$ est un morphisme fini étale, on dit encore que $X$ est un \emph{revêtement étale} de $Y$.

On notera que dans ce cas, $f$ est universellement ouvert et universellement fermé, et en particulier $f(X)$ est une partie à la fois ouverte et fermée de $Y$.

On dit qu'un revêtement étale $X$ de $Y$ est \emph{trivial} si $X$ est \emph{somme} d'un nombre fini de préschémas \emph{isomorphes} à $Y$. On dit qu'un revêtement étale $X$ de $Y$ est \emph{localement trivial} si le morphisme $f:X\to Y$ est tel que tout point $y\in Y$ ait un voisinage ouvert $U$ pour lequel le revêtement $f^{-1}(U)$ de $U$ soit trivial.

  \item[(ii)] Soit $f:X\to Y$ un morphisme fini, localement libre de rang $n$; posons $\sh{B}=f_*(\sh{O}_X)$ et soit $u=\operatorname{astr}_{\sh{B}|\sh{O}_Y}:\sh{B}\to\check{\sh{B}}$; on en déduit un homomorphisme puissance extérieure $n$-ème
\[
  \bigwedge^n u:\bigwedge^n\sh{B}\longrightarrow\bigwedge^n\check{\sh{B}}=(\bigwedge^n\sh{B})^\vee
\]
de $\sh{O}_Y$-Modules \emph{inversibles}, et par suite ($0_{\mathrm I}$, 5.4.2) un élément
\begin{equation}
  d_{X/Y}\in\Gamma\left(Y,(\bigwedge^n\check{\sh{B}})\otimes_{\sh{O}_Y}(\bigwedge^n\check{\sh{B}})\right)
  \tag{18.2.7.1}\label{IV.18.2.7.1}
\end{equation}
que l'on appelle le \emph{discriminant de $X$ sur $Y$}. D'ailleurs, comme $(\bigwedge^n\check{\sh{B}})\otimes_{\sh{O}_Y}(\bigwedge^n\check{\sh{B}})$ est le dual de $(\bigwedge^n\sh{B})\otimes_{\sh{O}_Y}(\bigwedge^n\sh{B})$, $d_{X/Y}$ peut aussi être identifié à un homomorphisme
\begin{equation}
  (\bigwedge^n\sh{B})\otimes_{\sh{O}_Y}(\bigwedge^n\sh{B})\longrightarrow\sh{O}_Y
  \tag{18.2.7.2}\label{IV.18.2.7.2}
\end{equation}
et on note $\mathscr{D}_{X/Y}$ l'Idéal quasi-cohérent de type fini de $\sh{O}_Y$, \emph{image} de l'homomorphisme (18.2.7.2), encore appelé \emph{Idéal discriminant de $X$ sur $Y$}.
\end{enumerate}
\end{remarks}

